\documentclass[11pt,a4paper]{article}
\usepackage[margin=19mm]{geometry}
\usepackage{amsmath,xcolor}
\usepackage{iftex}
\ifPDFTeX
  \usepackage{CJKutf8}
  % CJKutf8 supplies its own Unicode map for each Chinese subfont.
  \makeatletter
  \AddToHook{cmd/CJK@addcmap/before}{\pdfnobuiltintounicode\font@name}
  \makeatother
\else
  \usepackage[no-math]{fontspec}
  \IfFontExistsTF{Noto Serif CJK SC}{
    \setmainfont{Noto Serif CJK SC}
    \setsansfont{Noto Sans CJK SC}
  }{
    \setmainfont{FandolSong-Regular.otf}[BoldFont=FandolSong-Bold.otf]
    \setsansfont{FandolHei-Regular.otf}[BoldFont=FandolHei-Bold.otf]
  }
  \ifXeTeX
    \XeTeXlinebreaklocale "zh"
    \XeTeXlinebreakskip=0pt plus 0.5pt
  \fi
\fi
\pagestyle{empty}
\setlength{\parindent}{0pt}
\linespread{1.12}
\setlength{\parskip}{6pt}
\setlength{\abovedisplayskip}{7pt}
\setlength{\belowdisplayskip}{7pt}
\setlength{\abovedisplayshortskip}{3pt}
\setlength{\belowdisplayshortskip}{3pt}
\definecolor{accent}{HTML}{24545C}
\newcommand{\heading}[1]{\par\vspace{7pt}{\sffamily\bfseries\large\color{accent}#1}\par}
\begin{document}
\ifPDFTeX\begin{CJK*}{UTF8}{gbsn}\fi
{\sffamily\bfseries\LARGE 多个暴露对已发生结局的因果责任}\par\vspace{4pt}
用 $X_j\in\{0,1\}$ 表示第 $j$ 个候选原因对应的暴露，$Y=1$ 表示目标结局发生。
对一个已观测到 $E=\{X_1=\cdots=X_p=1,\,Y=1\}$ 的个体，我们问：这 $p$ 个暴露对这次结局各自承担多少责任？为此，需要比较同一个体在不同暴露组合下的潜在结果。

\heading{暴露组合与潜在结果}
记 $N=\{1,\ldots,p\}$，用 $\phi$ 表示空集。对任意 $S\subseteq N$，令 $Y(S)$ 表示仅将 $S$ 中的暴露设为 1、其余暴露均设为 0 时的潜在结果。例如，$Y(\phi)$ 是无暴露时的结局，$Y(N)$ 是所有暴露均存在时的结局。由一致性，在 $E$ 下有 $Y(N)=1$。

\heading{用 Shapley 值分配责任}
对同一个体，若仅有 $S$ 中的暴露，结局为 $Y(S)$。在此基础上再加入暴露 $j\notin S$，结局为 $Y(S\cup\{j\})$。因此，$j$ 在这一暴露组合下的边际贡献为 $Y(S\cup\{j\})-Y(S)$。
将 $p$ 个暴露的全部 $p!$ 种排列等权考虑，对 $j$ 的边际贡献取平均，得到暴露 $j$ 对该个体结局的责任
\[
 r_j=\sum_{S\subseteq N\setminus\{j\}}
 \frac{|S|!\,(p-|S|-1)!}{p!}\,\bigl[Y(S\cup\{j\})-Y(S)\bigr],
 \qquad j=1,\ldots,p .
\]
和式中的权重恰是随机排列中排在 $j$ 之前的暴露集合等于 $S$ 的概率。若该个体的潜在结果满足单调性，即对所有 $S$ 和 $j\notin S$ 均有 $Y(S\cup\{j\})\ge Y(S)$，则 $r_j$ 就是在均匀随机排列下，加入 $j$ 时结局从 0 变为 1 的概率。

定义背景责任 $r_0=Y(\phi)$。沿任一排列加总各步的边际贡献，均得到 $Y(N)-Y(\phi)$，所以在 $E$ 下有 $r_0+\sum_{j=1}^p r_j=1$。各人的责任取决于各自的潜在结果。对具有相同观测事实 $E$ 的人群，第 $j$ 项的平均责任定义为
\[
 R_j(E)=\mathrm{E}\!\left[r_j\mid X_1=\cdots=X_p=1,\,Y=1\right],\qquad j=0,1,\ldots,p .
\]

\heading{两个暴露与三个暴露的例子}
当 $p=2$ 时，将两种暴露分别设为 $a,b$ 时的潜在结果记为 $Y(a,b)$。$X_1$ 可以先加入，也可以在 $X_2$ 之后加入。两种顺序各占一半，因此
\[
 r_1=\tfrac12\bigl[Y(1,0)-Y(0,0)\bigr]
       +\tfrac12\bigl[Y(1,1)-Y(0,1)\bigr].
\]
当 $p=3$ 时，考虑这样一个个体：$X_1$ 单独足以使结局发生，或者 $X_2$ 与 $X_3$ 同时存在也足以使结局发生，其余组合下结局均为 0。在六种排列中，$X_1$ 排第一或第二时使结局从 0 变为 1，共四种。另两种排列分别由 $X_2$ 和 $X_3$ 触发结局。因此 $(r_1,r_2,r_3)=(2/3,1/6,1/6)$，背景责任为 0。

若原始结局 $O\in\{1,2,3\}$，且 3 级最严重，要解释“为何出现 3 级”，只需令 $Y(S)=\mathbf{1}\{O(S)=3\}$，即可沿用上述分配。

\vfill
{\footnotesize\color{accent} Shapley (1953), \textit{A Value for $n$-Person Games}. \quad
 Luo, Deng, Zhao \& Geng (2026), arXiv:2606.18459.}
\ifPDFTeX\end{CJK*}\fi
\end{document}
